% !PRINTED 715
\weberpage{715}

Bezeichnen wir daher mit
\[
(12)\quad E = \pm
\begin{vmatrix}
\log|\varepsilon_{1,1}|, & \log|\varepsilon_{1,2}|, & \ldots, & \log|\varepsilon_{1,\nu-1}| \\
\log|\varepsilon_{2,1}|, & \log|\varepsilon_{2,2}|, & \ldots, & \log|\varepsilon_{2,\nu-1}| \\
\cdots & \cdots & \cdots & \cdots \\
\log|\varepsilon_{\nu-1,1}|, & \log|\varepsilon_{\nu-1,2}|, & \ldots, & \log|\varepsilon_{\nu-1,\nu-1}|
\end{vmatrix}
\]
den Regulator des Körpers $H_m$ (\S.\,187), so ist bei der Anwendung des Ausdruckes (11) im Körper $\Omega_m$
\[
L = 2^{\nu-1} E,
\]
und im Körper $H_m$
\[
L = E
\]
zu setzen. Es ist ferner
\[
\begin{aligned}
\text{in } \Omega_m: & \quad n = \mu, \quad \nu = \tfrac{1}{2}\mu, \\
\text{in } H_m: & \quad n = \tfrac{1}{2}\mu, \quad \nu = \tfrac{1}{2}\mu.
\end{aligned}
\]
Im Körper $H_m$ sind nur die zwei Einheitswurzeln $\pm 1$ enthalten. In $\Omega_m$ sind die mit positivem und negativem Zeichen genommenen Potenzen von $r$, sonst aber keine Einheitswurzeln enthalten (\S.\,175), und bei der Zählung ist noch zu beachten, dass bei geradem $m$ unter den Potenzen von $r$ auch $-1$ enthalten ist, bei ungeradem $m$ nicht. Daher haben wir
\[
\begin{aligned}
\text{in } \Omega_m: & \quad w = 2m, \quad m \text{ ungerade} \\
& \quad = m, \quad m \text{ gerade,} \\
\text{in } H_m: & \quad w = 2.
\end{aligned}
\]
Endlich haben wir noch, wenn $\varDelta$, $\varDelta_1$ die Grundzahlen der Körper $\Omega_m$, $H_m$ sind (\S.\,176):
\[
\begin{aligned}
\pm \varDelta & = q\varDelta_1^2 \quad m \text{ ungerade} \\
& = 4\varDelta_1^2 \quad m \text{ gerade,}
\end{aligned}
\]
und es folgt, wenn wir der Einfachheit wegen $\nu$ für $\tfrac{1}{2}\mu$ setzen:
\[
\begin{aligned}
g & = \frac{2^{2(\nu-1)}\pi^\nu E}{m\sqrt{q}\,\varDelta_1} \quad m \text{ ungerade} \\
g & = \frac{2^{2(\nu-1)}\pi^\nu E}{m\varDelta_1} \quad m \text{ gerade,}
\end{aligned}
\]
\weq{13}{g_1 = \frac{2^{\nu-1}E}{\sqrt{\varDelta_1}},}
woraus
\weq{14}{\begin{aligned}
g & = \frac{2^{\nu-1}\pi^\nu g_1}{m\sqrt{q\varDelta_1}} \quad m \text{ ungerade} \\
& = \frac{2^{\nu-1}\pi^\nu g_1}{m\sqrt{\varDelta_1}} \quad m \text{ gerade.}
\end{aligned}}
